\chapter{Capitolo 13 - Benchmark, disegno sperimentale, replicabilità multi-dominio}

\section{1. Problema}

Il Capitolo 12 ha definito la grammatica metodologica della validazione OCT. Ma un protocollo, da solo, non garantisce prova robusta: serve un disegno benchmark che produca evidenza comparabile nel tempo, nei team e nei domini.

Nel lavoro ordinativo il rischio principale e doppio:

\begin{enumerate}
\def\labelenumi{\arabic{enumi}.}
\tightlist
\item
  benchmark troppo facili, che confermano qualsiasi ipotesi già favorita;
\item
  benchmark non comparabili, che impediscono di distinguere progresso reale da variazione accidentale.
\end{enumerate}

Per questo il problema del capitolo e preciso:

\textbf{come costruire un sistema benchmark multi-dominio che renda misurabile la differenza tra performance locale e validazione trasferibile, senza collassare la teoria in una pura ottimizzazione numerica?}

\subsection{1.1 Dal protocollo all'infrastruttura}

Il protocollo (Capitolo 12) risponde a ``come decidere''. Il benchmark design risponde a ``su cosa decidere''. La tenuta scientifica richiede entrambi.

Senza infrastruttura benchmark:

\begin{enumerate}
\def\labelenumi{\arabic{enumi}.}
\tightlist
\item
  i claim restano narrativi;
\item
  le repliche sono episodiche;
\item
  la decisione \texttt{validated} resta fragile.
\end{enumerate}

Con infrastruttura benchmark:

\begin{enumerate}
\def\labelenumi{\arabic{enumi}.}
\tightlist
\item
  i claim sono testati in condizioni comparabili;
\item
  le repliche indipendenti diventano obbligatorie, non opzionali;
\item
  la curva di evidenza può essere monitorata per versione.
\end{enumerate}

\subsection{\texorpdfstring{1.2 Scope della milestone \texttt{v1.4}}{1.2 Scope della milestone v1.4}}

Il Capitolo 13 ha funzione ponte tra metodo e governance:

\begin{enumerate}
\def\labelenumi{\arabic{enumi}.}
\tightlist
\item
  implementa la PV-8 in benchmark reali;
\item
  definisce la nozione operativa di indipendenza benchmark;
\item
  prepara il Capitolo 14 su audit trail, gate e versioning epistemico.
\end{enumerate}

\begin{center}\rule{0.5\linewidth}{0.5pt}\end{center}

\section{2. Tesi del capitolo}

La tesi e articolata in dieci enunciati.

\begin{enumerate}
\def\labelenumi{\arabic{enumi}.}
\tightlist
\item
  Un claim OCT e robusto solo se supera benchmark indipendenti per costruzione, non solo per nome.
\item
  L'indipendenza benchmark richiede separazione di dati, seed policy, tuning e log namespace.
\item
  La replicabilità multi-dominio e criterio strutturale, non supplemento statistico.
\item
  Le metriche canoniche (\texttt{Coh}, \texttt{Phi}, \texttt{Delta}) devono convivere con metriche specifiche per task.
\item
  La stratificazione dei casi e necessaria per evitare conferme trainate da sottogruppi facili.
\item
  L'analisi di sensibilità e parte del test principale, non appendice facoltativa.
\item
  Il fallimento parziale e informazione scientifica e va pubblicato con stessa disciplina dei successi.
\item
  Ogni benchmark deve includere policy di replay e criteri di rerun tracciati.
\item
  Un pacchetto benchmark senza manifest e hash non e considerato riproducibile oltre livello \texttt{R1}.
\item
  Il disegno proposto e sufficiente per sostenere la governance epistemica del Capitolo 14.
\end{enumerate}

\begin{center}\rule{0.5\linewidth}{0.5pt}\end{center}

\section{3. Definizioni canoniche}

\subsection{Definizione 13.1 (Benchmark indipendente)}

Due benchmark \texttt{B\_i}, \texttt{B\_j} sono indipendenti se valgono simultaneamente:

\begin{enumerate}
\def\labelenumi{\arabic{enumi}.}
\tightlist
\item
  \texttt{D\_i\ \ inter\ \ D\_j} trascurabile rispetto a soglia preregistrata;
\item
  seed policy distinta;
\item
  assenza di tuning trasversale post-freeze;
\item
  log namespace separato e auditabile;
\item
  shift di distribuzione dichiarato e verificabile.
\end{enumerate}

\subsection{Definizione 13.2 (Asse di dominio)}

Un asse di dominio \texttt{A\_k} e una famiglia coerente di contesti osservativi \texttt{Omega} che condividono:

\begin{enumerate}
\def\labelenumi{\arabic{enumi}.}
\tightlist
\item
  tipo di oggetti/morfismi osservati;
\item
  regime di rumore dominante;
\item
  criteri semantici di errore.
\end{enumerate}

Esempi operativi per OCT:

\begin{enumerate}
\def\labelenumi{\arabic{enumi}.}
\tightlist
\item
  \texttt{A\_struct}: strutture compositive e diagrammatiche;
\item
  \texttt{A\_pipeline}: processi multi-step di trasformazione;
\item
  \texttt{A\_social}: dinamiche narrative longitudinali.
\end{enumerate}

\subsection{Definizione 13.3 (Pacchetto benchmark minimo)}

\texttt{PB\_min\ :=\ \textless{}Manifest,\ Dataset,\ Script,\ Config,\ Hash,\ Report,\ DecisionTable\textgreater{}}.

Un benchmark e considerato pubblicabile solo se il pacchetto minimo e completo.

\subsection{Definizione 13.4 (Replicabilità multi-dominio)}

Un claim e multi-dominio replicabile se la conferma si osserva in almeno due assi \texttt{A\_k} con contesti non equivalenti banalmente.

\subsection{Definizione 13.5 (Stress strutturale)}

Stress strutturale: perturbazione controllata che aumenta difficolta compositiva mantenendo il significato del claim.

\subsection{Definizione 13.6 (Indice di trasferibilità)}

Per claim \texttt{C}:

\texttt{T\_C\ :=\ mean\_k(score\_k)\ -\ penalty\_var}

con:

\begin{enumerate}
\def\labelenumi{\arabic{enumi}.}
\tightlist
\item
  \texttt{score\_k} performance ordinativa su asse \texttt{A\_k};
\item
  \texttt{penalty\_var} penalità per variabilità eccessiva cross-dominio.
\end{enumerate}

\texttt{T\_C} alto indica robustezza trasferibile; \texttt{T\_C} basso indica risultato locale.

\begin{center}\rule{0.5\linewidth}{0.5pt}\end{center}

\section{4. Sviluppo formale}

\subsection{4.1 Architettura benchmark in tre livelli}

Proponiamo un'architettura a tre livelli.

\begin{enumerate}
\def\labelenumi{\arabic{enumi}.}
\tightlist
\item
  \textbf{Livello L1 (core test)}: verifica diretta del claim su benchmark principale.
\item
  \textbf{Livello L2 (indipendenza)}: replica su benchmark \#2 con shift esplicito.
\item
  \textbf{Livello L3 (stress/audit)}: sensibilità a soglie, ablation, replay indipendente.
\end{enumerate}

La regola e semplice:

\begin{enumerate}
\def\labelenumi{\arabic{enumi}.}
\tightlist
\item
  L1 senza L2 produce solo evidenza preliminare;
\item
  L1+L2 senza L3 produce evidenza promettente ma incompleta;
\item
  L1+L2+L3 e condizione standard per decisione forte.
\end{enumerate}

\subsection{4.2 Disegno degli assi di dominio}

Il framework OCT ha già evidenziato tre famiglie di test rilevanti:

\begin{enumerate}
\def\labelenumi{\arabic{enumi}.}
\tightlist
\item
  D02 (commutatività non produttiva) su asse strutturale;
\item
  A01 (stabilità pipeline IA) su asse operativo-procedurale;
\item
  D03 (perdita ontologica forgetting) su asse misto strutturale-operativo.
\end{enumerate}

Questa triade e utile perché copre:

\begin{enumerate}
\def\labelenumi{\arabic{enumi}.}
\tightlist
\item
  differenza puramente compositiva;
\item
  utilità pratica su pipeline reali;
\item
  perdita di struttura sotto trasformazioni canoniche.
\end{enumerate}

Il capitolo impone che ogni nuova campagna espliciti a priori quale asse copre e quale resta non coperto.

\subsection{4.3 Criteri operativi di indipendenza}

L'indipendenza non può essere dichiarata genericamente. Va verificata con checklist.

Checklist minima:

\begin{enumerate}
\def\labelenumi{\arabic{enumi}.}
\tightlist
\item
  \texttt{NO} riuso campioni dal benchmark precedente oltre soglia consentita;
\item
  split dati congelato prima del run;
\item
  separazione delle seed rispetto al ciclo precedente;
\item
  divieto di tuning post-freeze;
\item
  motivazione obbligatoria per rerun con hash log.
\end{enumerate}

Questi criteri sono coerenti con il manifest Cycle 3 e con la policy di execution lock già adottata.

\subsection{4.4 Stratificazione campionaria e copertura}

Ogni benchmark deve esplicitare stratificazione su tre assi:

\begin{enumerate}
\def\labelenumi{\arabic{enumi}.}
\tightlist
\item
  complessità strutturale;
\item
  livello di rumore;
\item
  ambiguità contestuale.
\end{enumerate}

Se uno strato e sottocampionato, il claim non può essere generalizzato. In pratica:

\begin{enumerate}
\def\labelenumi{\arabic{enumi}.}
\tightlist
\item
  risultati aggregati senza breakdown di strato sono considerati incompleti;
\item
  la decisione può al massimo restare in \texttt{revise}.
\end{enumerate}

\subsection{4.5 Piano statistico minimo}

Senza imporre un unico paradigma statistico, fissiamo cinque obblighi minimi:

\begin{enumerate}
\def\labelenumi{\arabic{enumi}.}
\tightlist
\item
  ipotesi nulla/alternativa dichiarate;
\item
  soglia \texttt{alpha} preregistrata;
\item
  misura di effetto oltre al p-value;
\item
  intervalli di confidenza o stima di incertezza equivalente;
\item
  criterio di successo aggregato multi-contesto.
\end{enumerate}

Il capitolo privilegia inferenza trasparente rispetto a ottimizzazione retorica di singole metriche.

\subsection{4.6 Potenza e numerosità campionaria}

La numerosità non può essere scelta solo per convenienza computazionale. Deve essere giustificata da:

\begin{enumerate}
\def\labelenumi{\arabic{enumi}.}
\tightlist
\item
  variabilità attesa;
\item
  effetto minimo di interesse;
\item
  livello di confidenza richiesto per decisione.
\end{enumerate}

Se la potenza e insufficiente, un esito ``non significativo'' non autorizza rigetto del claim; autorizza solo revisione del disegno.

\subsection{4.7 Multi-contesto e regola 2/3}

In continuità con i cicli già eseguiti, adottiamo una regola pragmatica:

\begin{enumerate}
\def\labelenumi{\arabic{enumi}.}
\tightlist
\item
  conferma minima in almeno \texttt{2/3} contesti preregistrati;
\item
  il terzo contesto può fallire solo con analisi causale esplicita.
\end{enumerate}

La regola evita due estremi:

\begin{enumerate}
\def\labelenumi{\arabic{enumi}.}
\tightlist
\item
  rigidità eccessiva (richiedere successo perfetto in ogni micro-variante);
\item
  permissività eccessiva (accettare successo in un solo contesto favorevole).
\end{enumerate}

\subsection{4.8 Instrumentation e audit trail}

La replicabilità dipende più dai log che dalle intenzioni. Ogni run deve produrre:

\begin{enumerate}
\def\labelenumi{\arabic{enumi}.}
\tightlist
\item
  manifest versione dati e script;
\item
  checksum artefatti;
\item
  parametri runtime;
\item
  output metrici raw e aggregati;
\item
  decision trace leggibile.
\end{enumerate}

Senza questa catena, il risultato resta al massimo \texttt{R1} e non supporta validazione forte.

\subsection{4.9 Tassonomia degli errori}

Per evitare diagnosi vaghe introduciamo una tassonomia minima degli errori benchmark:

\begin{enumerate}
\def\labelenumi{\arabic{enumi}.}
\tightlist
\item
  \texttt{E\_data}: incoerenza o leakage nei dati;
\item
  \texttt{E\_model}: instabilità del mapping o della pipeline;
\item
  \texttt{E\_metric}: mismatch tra metrica e significato del claim;
\item
  \texttt{E\_context}: trasferimento improprio tra domini;
\item
  \texttt{E\_decision}: applicazione incoerente della matrice decisionale.
\end{enumerate}

Ogni errore va classificato in severità (\texttt{low}, \texttt{medium}, \texttt{high}) e associato ad azione correttiva.

\subsection{4.10 Policy di rerun e freezing}

Il rerun e lecito solo se:

\begin{enumerate}
\def\labelenumi{\arabic{enumi}.}
\tightlist
\item
  motivazione tecnica documentata;
\item
  stessa preregistrazione o nuova preregistrazione versionata;
\item
  dichiarazione esplicita di cosa cambia e cosa resta invariato.
\end{enumerate}

Il freezing conclude una fase sperimentale e blocca la retro-modifica opportunistica dei criteri.

\subsection{4.11 Disegno per pubblicazione aperta}

Un benchmark scientificamente utile e anche pubblicabile in forma verificabile.

Pacchetto raccomandato:

\begin{enumerate}
\def\labelenumi{\arabic{enumi}.}
\tightlist
\item
  repository operativo (GitHub) con script e manifest;
\item
  snapshot DOI (Zenodo) della release validata;
\item
  workspace progettuale (OSF) con preregistrazioni e materiale intermedio;
\item
  eventuale mirror dataset (Hugging Face) con dataset card metodologica.
\end{enumerate}

La chiave non e moltiplicare piattaforme, ma preservare coerenza di versione tra piattaforme.

\subsection{4.12 Condizione di sufficienza del capitolo}

Il capitolo e metodologicamente sufficiente se produce tre risultati:

\begin{enumerate}
\def\labelenumi{\arabic{enumi}.}
\tightlist
\item
  definizione operativa di benchmark indipendente;
\item
  schema di replicabilità multi-dominio con criteri espliciti;
\item
  pipeline pubblicabile e auditabile end-to-end.
\end{enumerate}

Questa condizione e soddisfatta dalle definizioni 13.1-13.6 e dai blocchi 4.1-4.11.

\subsection{\texorpdfstring{4.13 Indice di qualità benchmark \texttt{Q\_B}}{4.13 Indice di qualità benchmark Q\_B}}

Per confrontare campagne diverse introduciamo un indicatore sintetico:

\texttt{Q\_B\ :=\ a*Ind\ +\ b*Rep\ +\ c*Trace\ +\ d*Stress\ -\ e*Leak}

dove:

\begin{enumerate}
\def\labelenumi{\arabic{enumi}.}
\tightlist
\item
  \texttt{Ind} misura l'indipendenza effettiva tra benchmark;
\item
  \texttt{Rep} misura il tasso di repliche riuscite;
\item
  \texttt{Trace} misura la completezza dell'audit trail;
\item
  \texttt{Stress} misura copertura e severità dei test di sensibilità;
\item
  \texttt{Leak} penalizza leakage dati e tuning opportunistico.
\end{enumerate}

Vincoli operativi:

\begin{enumerate}
\def\labelenumi{\arabic{enumi}.}
\tightlist
\item
  coefficienti \texttt{a,b,c,d,e\ \textgreater{}=\ 0};
\item
  \texttt{a+b+c+d+e\ =\ 1};
\item
  soglia minima preregistrata per rilascio esterno.
\end{enumerate}

\texttt{Q\_B} non sostituisce la decisione teorica sul claim, ma decide la prontezza metodologica del pacchetto benchmark.

\subsection{4.14 Protocollo di replica indipendente}

Per dichiarare una replica realmente indipendente proponiamo il seguente protocollo minimo:

\begin{enumerate}
\def\labelenumi{\arabic{enumi}.}
\tightlist
\item
  team esecutore differente da quello del benchmark originario;
\item
  ambiente ricostruito da zero usando solo artefatti pubblici;
\item
  verifica hash di dataset/script prima del run;
\item
  esecuzione senza modifica dei criteri decisionali preregistrati;
\item
  report finale con confronto puntuale tra esito originario e replica.
\end{enumerate}

La replica indipendente non e una formalità: e il passaggio che trasforma un risultato interno in conoscenza affidabile per la comunità.

\begin{center}\rule{0.5\linewidth}{0.5pt}\end{center}

\section{5. Proposizioni e teoremi locali}

\subsection{Proposizione 13.1 (Insufficienza del benchmark nominalmente diverso)}

Enunciato:
Due benchmark con nomi differenti ma senza separazione di dati/seed/tuning non forniscono evidenza indipendente.

Intuizione:
l'indipendenza e proprietà del processo, non dell'etichetta.

Stato: \texttt{validated}.

\subsection{Proposizione 13.2 (Necessità della stratificazione)}

Enunciato:
Senza stratificazione esplicita, la performance aggregata può sovrastimare la robustezza di un claim.

Intuizione:
la media può nascondere collassi locali ad alta rilevanza teorica.

Stato: \texttt{validated}.

\subsection{Proposizione 13.3 (Valore informativo del fallimento)}

Enunciato:
Nel contesto OCT, un fallimento tracciato e classificato aumenta informazione metodologica più di un successo non auditabile.

Intuizione:
il fallimento ben localizzato orienta revisione teorica e disegno successivo.

Stato: \texttt{in\_review}.

\subsection{Teorema 13.4 (Sufficienza operativa del disegno L1-L2-L3)}

Ipotesi:

\begin{enumerate}
\def\labelenumi{\arabic{enumi}.}
\tightlist
\item
  benchmark L1 con preregistrazione completa;
\item
  benchmark L2 indipendente secondo Def. 13.1;
\item
  stress/audit L3 con report riproducibile.
\end{enumerate}

Tesi:
Il claim testato dispone di evidenza metodologicamente sufficiente per decisione \texttt{validated/revise/reject} non arbitraria.

Proof sketch:

\begin{enumerate}
\def\labelenumi{\arabic{enumi}.}
\tightlist
\item
  L1 fornisce test diretto del claim;
\item
  L2 limita overfitting al benchmark iniziale;
\item
  L3 misura robustezza a perturbazioni e integrità della catena di prova.
\end{enumerate}

Conseguenze:
fornisce standard minimo per promozione di stato in catalogo teorematico OCT.

Stato: \texttt{in\_proof}.

\subsection{Teorema 13.5 (Necessità della replicabilità multi-dominio)}

Ipotesi:

\begin{enumerate}
\def\labelenumi{\arabic{enumi}.}
\tightlist
\item
  claim confermato solo su un asse di dominio;
\item
  assenza di repliche su asse indipendente;
\item
  elevata variabilità potenziale tra contesti.
\end{enumerate}

Tesi:
La conferma mono-dominio non e sufficiente per inferenza generale in OCT.

Proof sketch:

\begin{enumerate}
\def\labelenumi{\arabic{enumi}.}
\tightlist
\item
  la semantica ordinativa dipende da \texttt{Omega};
\item
  cambi di dominio alterano distribuzioni e vincoli compositivi;
\item
  senza replica multi-asse non si distingue effetto teorico da compatibilità locale.
\end{enumerate}

Conseguenze:
legittima la regola di replicazione minima su assi multipli.

Stato: \texttt{in\_proof}.

\begin{center}\rule{0.5\linewidth}{0.5pt}\end{center}

\section{6. Implicazioni operative}

\subsection{6.1 Per il Capitolo 14 (governance epistemica)}

Il capitolo consegna a \texttt{v1.4} tre oggetti che il Capitolo 14 dovra governare:

\begin{enumerate}
\def\labelenumi{\arabic{enumi}.}
\tightlist
\item
  versioni benchmark e relativi freeze;
\item
  audit trail dei rerun;
\item
  criteri di accettazione legati a livelli di riproducibilità.
\end{enumerate}

\subsection{6.2 Per il programma di validazione teoremi}

La struttura L1-L2-L3 permette di ordinare i teoremi in priorità:

\begin{enumerate}
\def\labelenumi{\arabic{enumi}.}
\tightlist
\item
  teoremi ad alto impatto con percorso completo immediato;
\item
  teoremi esplorativi con L1 rapido e ingresso in coda L2/L3;
\item
  teoremi in conflitto con focus su stress e ablation prima del rilascio.
\end{enumerate}

\subsection{6.3 Per pubblicazione e revisione esterna}

Un revisore esterno deve poter rispondere a tre domande senza ambiguità:

\begin{enumerate}
\def\labelenumi{\arabic{enumi}.}
\tightlist
\item
  che cosa e stato testato;
\item
  come e stato testato;
\item
  perché la decisione finale e giustificata.
\end{enumerate}

Il capitolo offre una struttura che rende queste risposte ispezionabili.

\subsection{6.4 Per roadmap 2026-2030}

La replicabilità multi-dominio e il criterio che separa:

\begin{enumerate}
\def\labelenumi{\arabic{enumi}.}
\tightlist
\item
  risultati locali da risultati fondativi;
\item
  prototipi interni da standard scientifici pubblicabili.
\end{enumerate}

Quindi la roadmap non dovrebbe contare solo ``numero di teoremi testati'', ma ``numero di teoremi con L1-L2-L3 completo''.

\subsection{6.5 Per onboarding di nuovi laboratori}

Un laboratorio esterno può entrare nel programma OCT senza conoscere tutta la storia del progetto, se riceve un kit standardizzato:

\begin{enumerate}
\def\labelenumi{\arabic{enumi}.}
\tightlist
\item
  guida iniziale con lessico minimo condiviso;
\item
  pacchetto \texttt{PB\_min} di un claim campione;
\item
  checklist di replica indipendente;
\item
  template report con sezione errori obbligatoria.
\end{enumerate}

Questo riduce la barriera d'ingresso e aumenta la probabilità di validazioni indipendenti reali.

\begin{center}\rule{0.5\linewidth}{0.5pt}\end{center}

\section{7. Limiti e non-claim}

Limiti espliciti:

\begin{enumerate}
\def\labelenumi{\arabic{enumi}.}
\tightlist
\item
  il capitolo definisce una baseline metodologica, non l'unica possibile;
\item
  alcuni domini richiedono metriche ausiliarie altamente specializzate;
\item
  la regola 2/3 e pragmatica, non dogmatica, e può richiedere adattamento motivato.
\end{enumerate}

Failure mode riconosciuti:

\begin{enumerate}
\def\labelenumi{\arabic{enumi}.}
\tightlist
\item
  benchmark formalmente indipendenti ma semanticamente ridondanti;
\item
  stress test costruiti troppo deboli per mettere in crisi il claim;
\item
  rerun usati per ``cercare'' un esito desiderato;
\item
  scarsa documentazione dei casi falliti.
\end{enumerate}

Contromisure raccomandate:

\begin{enumerate}
\def\labelenumi{\arabic{enumi}.}
\tightlist
\item
  revisione indipendente del manifest prima dell'esecuzione;
\item
  quota minima di casi ``difficili'' in ogni strato;
\item
  obbligo di pubblicazione dei fallimenti di severità \texttt{high};
\item
  audit periodico su coerenza tra preregistrazione e report finale.
\end{enumerate}

Box C - Non-claim

Questo capitolo non implica:

\begin{enumerate}
\def\labelenumi{\arabic{enumi}.}
\tightlist
\item
  che ogni claim debba superare tutti i domini possibili;
\item
  che la significatività statistica da sola determini la validità teorica;
\item
  che un benchmark ben progettato sostituisca la necessità di formalizzazione matematica.
\end{enumerate}

\begin{center}\rule{0.5\linewidth}{0.5pt}\end{center}

\section{8. Claim Ledger (S0-S3)}

{\def\LTcaptype{none} % do not increment counter
\begin{longtable}[]{@{}lllll@{}}
\toprule\noalign{}
ID & Claim & Tipo & Evidenza & Stato \\
\midrule\noalign{}
\endhead
\bottomrule\noalign{}
\endlastfoot
C13.1 & L'indipendenza benchmark richiede separazione di dati, seed, tuning e log & metodologico & S1 & validated \\
C13.2 & La stratificazione e necessaria per evitare sovrastima della robustezza & statistico-metodologico & S1 & validated \\
C13.3 & Il disegno L1-L2-L3 e sufficiente come baseline per decisioni non arbitrarie & metodologico-formale & S2 & in\_proof \\
C13.4 & La replicabilità multi-dominio e condizione necessaria per inferenza generale OCT & epistemico & S2 & in\_proof \\
C13.5 & La pubblicazione dei fallimenti migliora la qualità cumulativa del programma & operativo & S1 & in\_review \\
C13.6 & Il pacchetto benchmark minimo \texttt{PB\_min} supporta rilascio multi-piattaforma coerente & infrastrutturale & S1 & in\_review \\
\end{longtable}
}

\begin{center}\rule{0.5\linewidth}{0.5pt}\end{center}

\section{9. Bridge al Capitolo 14}

Il Capitolo 13 ha stabilito come progettare benchmark indipendenti e replicabili, e come trasformare il protocollo in infrastruttura sperimentale.

Il Capitolo 14 completera la milestone \texttt{v1.4} introducendo la governance epistemica:

\begin{enumerate}
\def\labelenumi{\arabic{enumi}.}
\tightlist
\item
  regole di versioning dei claim;
\item
  audit trail decisionale;
\item
  quality gate per release pubbliche;
\item
  politica di manutenzione del corpus teorematico.
\end{enumerate}

In breve: il Capitolo 13 costruisce la macchina sperimentale; il Capitolo 14 ne definisce la costituzione operativa.

\begin{center}\rule{0.5\linewidth}{0.5pt}\end{center}

\section{Riferimenti}

Riferimenti di framework interno:

\begin{enumerate}
\def\labelenumi{\arabic{enumi}.}
\tightlist
\item
  Ghioni, F. (2026). \texttt{TE\_CORE\_v5.1.md}.
\item
  Ghioni, F. (2026). \texttt{Ordinative\_Set\_Theory\_OST\_A\_Concise\_Guide\_For\_AI\_v2\_1.md}.
\item
  \texttt{OCT\_VALIDATION\_PROTOCOL\_v0\_1.md}.
\item
  \texttt{OCT\_Theory\_and\_Theorems/validation/CYCLE\_3\_2026-04-18/BENCHMARK\_MANIFEST\_cycle3\_v0\_1.md}.
\item
  \texttt{OCT\_Theory\_and\_Theorems/validation/CYCLE\_4\_2026-04-19/CYCLE\_4\_REPORT.md}.
\item
  \texttt{OCT\_FOUNDATIONAL\_BOOK\_CHAPTER\_12\_v1\_0.md}.
\end{enumerate}

Riferimenti primari:

\begin{enumerate}
\def\labelenumi{\arabic{enumi}.}
\tightlist
\item
  Popper, K. (1959). \emph{The Logic of Scientific Discovery}.
\item
  Lakatos, I. (1978). \emph{The Methodology of Scientific Research Programmes}.
\item
  Nosek, B. et al. (2018). \emph{The preregistration revolution}.
\item
  Goodman, S., Fanelli, D., Ioannidis, J. (2016). \emph{What does research reproducibility mean?}.
\end{enumerate}
